% =====================================================================
\section{Model 1: Physics, Estimation and Geometric Control}
\label{sec:model1}
% =====================================================================

Model~1 is the loop that flies the drone (\cref{fig:flow1}). Blocks 1--3 describe how the drone moves and how a small
error grows. Block 4 estimates the state from the sensors. Blocks 5--7 decide thrust and torque. Block 8 turns them
into motor speeds. Toy model files: \code{M1\_1\_Dynamics\_Linearization.mlx} (blocks 1--3),
\code{M1\_2\_Error\_State\_EKF.mlx} (block 4) and\\ \code{M1\_3\_SO3\_Control\_Closed\_Loop.mlx} (blocks 5--8).

% ---------------------------------------------------------------------
\subsection{Block 1 --- How the drone moves (nonlinear dynamics)}
\label{sec:m1b1}
% ---------------------------------------------------------------------

\begin{idea}
A quadrotor is one solid body. Its four propellers can only push in one direction: straight out of the top of
the drone. Gravity pulls it down. To move sideways it must first \emph{tilt}, so that part of its push points
sideways. To tilt, the rotors push unevenly, which creates a twisting force (torque). So there are two
connected stories: how the body \emph{moves} (Newton) and how it \emph{turns} (Euler).
\end{idea}

\begin{eqbox}{Equations (1)--(4): rigid-body dynamics of the quadrotor}
\begin{align}
  \dot{\vp} &= \vv \tag{1}\label{eq:1}\\
  \dot{\vv} &= \vg + \frac{T}{m}\,R\,\ve_3 \tag{2}\label{eq:2}\\
  \dot{R}   &= R\,\skw{\vw} \tag{3}\label{eq:3}\\
  \dot{\vw} &= J^{-1}\big(\vtau - \vw\times J\vw\big) \tag{4}\label{eq:4}
\end{align}
State $\vx=(\vp,\vv,R,\vw)$, input $\vu=[T;\ \vtau]$, $\ve_3=[0,0,1]^\top$. In short: $\dot\vx=\vf(\vx,\vu)$.
\end{eqbox}

\subsubsection*{Where each equation comes from}

\textbf{(1) Position.} Velocity is, by definition, how fast position changes: $\vv=d\vp/dt$. Nothing more.

\textbf{(2) Newton's second law, $m\,\dot{\vv}=$ sum of forces.} Two forces act on the drone.
\begin{enumerate}
  \item Gravity: $m\vg=m[0,0,-9.81]^\top$, always straight down.
  \item Thrust: the four rotors together push with $T$ newtons along the drone's own up-axis. In the body frame this
        is $T\ve_3=[0,0,T]^\top$. To write it in the world frame we multiply by $R$ (\cref{sec:rotations}):
        $R\,T\ve_3=T\,\vb_3$, where $\vb_3$ is the third column of $R$, the direction the top of the drone points.
\end{enumerate}
So $m\dot\vv=m\vg+T R\ve_3$. Dividing by $m$ gives (2). Air drag and wind are left out of the model. In the toy
simulation they are added as unknown ``gusts'', which is exactly what the controller has to cope with.

\textbf{(3) How the attitude changes.} Take a point fixed on the drone, at $\bm r$ in body coordinates. In the world
it is at $R\bm r$, so its velocity is $\dot R\,\bm r$. A body spinning with angular velocity $\vw$ (measured in the body)
moves that point with velocity $\vw\times\bm r$ in body coordinates, which is $R(\vw\times\bm r)=R\skw{\vw}\bm r$ in
the world. Since this holds for every $\bm r$: $\dot R=R\skw{\vw}$. This is how the gyroscope reading $\vw$ turns the
attitude.

\textbf{(4) Euler's law for rotation.} The ``spin momentum'' (angular momentum) is $J\vw$, where $J$ is the
inertia matrix (how hard the body is to spin about each axis). Torque changes angular momentum. Because the body
frame itself is turning, an extra term appears:
\[
J\dot{\vw}+\vw\times J\vw=\vtau\quad\Longrightarrow\quad \dot\vw=J^{-1}(\vtau-\vw\times J\vw).
\]
The term $\vw\times J\vw$ is the \textbf{gyroscopic coupling}: spinning about two axes at once starts a spin about the
third axis, even with no torque. (Derivation: the world-frame momentum is $\bm h=RJ\vw$; Newton for rotation says
$\dot{\bm h}=R\vtau$; differentiating with (3) gives $R(\skw{\vw}J\vw+J\dot\vw)=R\vtau$, which is the equation above.)

\begin{toy}[hover, a 10 degree tilt, and gyroscopic coupling ($m=1.5$\,kg, $J=\diag(0.02,0.02,0.04)$)]
\textbf{1. Hover.} Level ($R=I$), not spinning, thrust equal to the weight $T=mg_0=1.5\times9.81=14.715$\,N:
\[
\dot\vv=\begin{bmatrix}0\\0\\-9.81\end{bmatrix}+\frac{14.715}{1.5}\begin{bmatrix}0\\0\\1\end{bmatrix}=\begin{bmatrix}0\\0\\0\end{bmatrix}.
\]
No acceleration: the drone hangs still. This balance point is called the \textbf{hover equilibrium}.

\textbf{2. Rolled by 10$^\circ$, same thrust.} The top of the drone now points along
$R\ve_3=[0,\ -\sin10^\circ,\ \cos10^\circ]=[0,\ -0.1736,\ 0.9848]$:
\[
\dot\vv=\begin{bmatrix}0\\0\\-9.81\end{bmatrix}+9.81\begin{bmatrix}0\\-0.1736\\0.9848\end{bmatrix}=\begin{bmatrix}0\\-1.7035\\-0.1490\end{bmatrix}\text{ m/s}^2 .
\]
It accelerates sideways (this is how a drone moves sideways!) and also sinks a little, because only 98.5\% of the
thrust still holds it up.

\textbf{3. Gyroscopic coupling.} Spinning at $\vw=[1,0,1]$\,rad/s with no torque:
$J\vw=[0.02,\,0,\,0.04]$, $\vw\times J\vw=[0\cdot0.04-1\cdot0,\ 1\cdot0.02-1\cdot0.04,\ 0]=[0,\,-0.02,\,0]$, so
$\dot\vw=-J^{-1}[0,-0.02,0]=[0,\ 1,\ 0]$\,rad/s$^2$. A spin about $y$ appears on its own.
\end{toy}

\begin{toycode}{M1\_1\_Dynamics\_Linearization}{Block 1}
x_hover = DroneLib.pack([0;0;2], [0;0;0], eye(3), [0;0;0]);
u_hover = [m*P.g0; 0; 0; 0]          % [T; tau_x; tau_y; tau_z]
xdot  = DroneLib.f(x_hover, u_hover, P);
v_dot = xdot(4:6)'                   % zero : the drone hangs still
\end{toycode}
The function \code{DroneLib.f} is equations (1)--(4) written in MATLAB, line by line:
\begin{lstlisting}[style=mat, title={\small\sffamily toy\_model3/DroneLib.m, function f}]
v_dot = P.g + u(1)/P.m*R*P.e3 + d(1:3)/P.m;          % (2) + wind force d
R_dot = R*DroneLib.hat(w);                            % (3)
w_dot = P.J \ (u(2:4) - cross(w, P.J*w) + d(4:6));    % (4) + disturbance torque
xdot  = [v; v_dot; R_dot(:); w_dot];                  % (1): p_dot = v
\end{lstlisting}

\textbf{A drone cannot fly by itself.} The toy model gives the hovering drone a tiny roll torque of 0.01\,N\,m for
just 0.1\,s and then lets it go with hover thrust (solved with \code{ode45}, as in the course notes). Nobody corrects
it, and after 4\,s it is \textbf{4.93\,m} away from where it started (\cref{fig:openloop}). A quadrotor is
\emph{unstable}: it needs to know its state (block 4) and it needs feedback (blocks 5--8).

\begin{figure}[H]
  \centering
  \includegraphics[width=0.62\linewidth]{M1_1_Dynamics_Linearization_1.png}
  \caption{Open loop: a 0.1\,s torque pulse and no correction. The drone slides further and further away.}
  \label{fig:openloop}
\end{figure}

\begin{real}
Equations (1)--(4) are ``the real drone''. In our setup Gazebo solves them for the Iris quadrotor; we never solve them
in the controller. Our node only \emph{reads} their result. The position and the attitude arrive on
\code{/ap/v1/pose/filtered}; the attitude comes as a quaternion and is turned into the rotation matrix $R$:
\pyfile{371}{386}{\node}{drones\_controller/controller\_node.py --- pose\_callback()}
\pyfile{24}{40}{\sad}{drones\_controller/state\_adapter.py --- quaternion\_to\_rotation\_matrix(): the matrix $R$}
The node also checks that this really is a rotation matrix ($R^\top R=I$, $\det R=1$), the two properties of
\cref{sec:rotations}:
\pyfile{127}{138}{\sad}{drones\_controller/state\_adapter.py --- validate\_rotation\_matrix()}
\end{real}

% ---------------------------------------------------------------------
\subsection{Block 2 --- Linearization: zooming in on the dynamics}
\label{sec:m1b2}
% ---------------------------------------------------------------------

\begin{idea}
Equations (1)--(4) are curved (nonlinear): they contain $R$, products and sines. But if we zoom in close to one
flight condition, any smooth curve looks like a straight line (\cref{sec:taylor}). Near that condition the drone
behaves like a simple \emph{linear} system, and linear systems are easy to predict and to control. The ``slopes'' of
this straight-line model are two matrices, $A$ and $B$.
\end{idea}

\begin{eqbox}{Equations (5)--(6): first order Taylor expansion and the Jacobians}
\begin{align}
  \vf(\vx,\vu) &\approx \vf(\hx,\hat\vu) + A\,\delta\vx + B\,\delta\vu \tag{5}\label{eq:5}\\
  A &= \left.\frac{\partial\vf}{\partial\vx}\right|_{\hx,\hat\vu}\in\R^{12\times12},\qquad
  B = \left.\frac{\partial\vf}{\partial\vu}\right|_{\hx,\hat\vu}\in\R^{12\times4} \tag{6}\label{eq:6}
\end{align}
with the 12-number error $\delta\vx=[\delta\vp;\ \delta\vv;\ \delta\vth;\ \delta\vw]$, where $R=\hR\,\Exp(\delta\vth)$, and $\delta\vu=\vu-\hat\vu$.
\end{eqbox}

\subsubsection*{What is $\delta\vx$?}
$\delta\vx$ is ``the difference between the true state and a reference state'' (later: the estimate).
For position, velocity and angular velocity this is a simple subtraction: $\delta\vp=\vp-\hp$ and so on. For the
attitude we cannot subtract matrices (the result would not be a rotation). Instead we ask: ``which small rotation
$\delta\vth$ takes $\hR$ to $R$?'', i.e.\ $R=\hR\,\Exp(\delta\vth)$, so $\delta\vth=\Log(\hR^\top R)$. That is 3
numbers. So $\delta\vx$ has $3+3+3+3=12$ numbers.

\subsubsection*{Deriving $A$ and $B$, one row at a time}
Put $\vp=\hp+\delta\vp$, $\vv=\hv+\delta\vv$, $R\approx\hR(I+\skw{\delta\vth})$, $\vw=\hw+\delta\vw$,
$T=\hat T+\delta T$, $\vtau=\hat\vtau+\delta\vtau$ into (1)--(4). Keep every term that contains \emph{one} small
quantity and throw away products of two small quantities (they are tiny squared).

\textbf{Position row.} $\dot\vp=\vv$ gives $\delta\dot{\vp}=\delta\vv$ directly.

\textbf{Velocity row.}
$\dot{\vv}\approx\vg+\frac{\hat T+\delta T}{m}\hR(I+\skw{\delta\vth})\ve_3$. Multiply out and subtract the reference
part $\vg+\frac{\hat T}{m}\hR\ve_3$:
\[
\delta\dot\vv=\frac{\hat T}{m}\hR\skw{\delta\vth}\ve_3+\frac{\delta T}{m}\hR\ve_3+\underbrace{\tfrac{\delta T}{m}\hR\skw{\delta\vth}\ve_3}_{\text{two small terms: drop}} .
\]
Using $\skw{\delta\vth}\ve_3=\delta\vth\times\ve_3=-\ve_3\times\delta\vth=-\skw{\ve_3}\delta\vth$:
\[
\delta\dot{\vv}=-\frac{\hat T}{m}\hR\skw{\ve_3}\,\delta\vth+\frac{1}{m}\hR\ve_3\,\delta T .
\]

\textbf{Attitude row.} Both $R$ and $\hR$ follow (3). Differentiating $R=\hR\Exp(\delta\vth)$ and keeping first order
terms gives
\[
\delta\dot{\vth}=-\skw{\hw}\,\delta\vth+\delta\vw .
\]
In words: an angular-velocity error makes the attitude error grow; the first term just accounts for the reference
frame itself spinning.

\textbf{Angular-velocity row.} The only nonlinear part of (4) is $\vw\times J\vw$. Its derivative with respect to $\vw$
is $\skw{\vw}J-\skw{J\vw}$ (product rule for the cross product), so
\[
\delta\dot{\vw}=J^{-1}\big(\skw{J\hw}-\skw{\hw}J\big)\delta\vw+J^{-1}\delta\vtau .
\]

Stacking the four rows (each block is $3\times3$; the columns follow $\delta\vp,\delta\vv,\delta\vth,\delta\vw$ for $A$
and $\delta T,\delta\vtau$ for $B$):
\begin{equation}
\label{eq:AB}\tag{6a}
A=\begin{bmatrix}
0&I&0&0\\
0&0&-\frac{\hat T}{m}\hR\skw{\ve_3}&0\\
0&0&-\skw{\hw}&I\\
0&0&0&J^{-1}\big(\skw{J\hw}-\skw{\hw}J\big)
\end{bmatrix},\qquad
B=\begin{bmatrix}0&0\\ \frac1m\hR\ve_3&0\\ 0&0\\ 0&J^{-1}\end{bmatrix}.
\end{equation}

\begin{toy}[the Jacobians at hover, and how good ``first order'' is]
At hover ($\hR=I$, $\hw=0$, $\hat T=mg_0$) the important blocks are
\[
-g_0\skw{\ve_3}=\begin{bmatrix}0&9.81&0\\-9.81&0&0\\0&0&0\end{bmatrix},\qquad
\tfrac1m\ve_3=\begin{bmatrix}0\\0\\0.667\end{bmatrix},\qquad
J^{-1}=\diag(50,\,50,\,25).
\]
Read them as sentences:
\begin{itemize}
  \item a pitch error of $\delta\theta_y$ radians accelerates the drone along $+x$ at $9.81\,\delta\theta_y$\,m/s$^2$
        (row 1 of the first matrix); a roll error accelerates it along $-y$ (row 2);
  \item 1\,N of extra thrust gives 0.667\,m/s$^2$ of extra \emph{vertical} acceleration and nothing else;
  \item 1\,N\,m of roll torque gives 50\,rad/s$^2$ of roll acceleration.
\end{itemize}
\textbf{Check by numbers.} The toy model also computes the derivatives numerically, $(\vf(\vx+h)-\vf(\vx-h))/2h$,
at a tilted, spinning, climbing state. The largest difference from the formulas is $4.9\times10^{-10}$ for $A$ and
$6.8\times10^{-10}$ for $B$: the derivation is right.

\textbf{How far can we trust the straight line?} Pitch the hovering drone by $\theta$. The linear model says
$\dot v_x=g_0\theta$, the truth is $g_0\sin\theta$ (the $\sin$ example of \cref{sec:taylor}):
\begin{center}\small
\begin{tabular}{rrrr}
\toprule
$\theta$ [deg] & linear $g_0\theta$ [m/s$^2$] & true $g_0\sin\theta$ [m/s$^2$] & error\\
\midrule
1 & 0.1712 & 0.1712 & 0.005\%\\
2.87 & 0.4914 & 0.4912 & 0.04\%\\
10 & 1.7122 & 1.7035 & 0.51\%\\
20 & 3.4243 & 3.3552 & 2.06\%\\
28.65 & 4.9054 & 4.7035 & 4.29\%\\
\bottomrule
\end{tabular}
\end{center}
In normal flight the tilt stays below 10$^\circ$, where the linear model is better than 0.5\%.
\end{toy}

\begin{toycode}{M1\_1\_Dynamics\_Linearization}{Block 2}
[A, B] = DroneLib.jacobians(x_hover, u_hover, P);
A_v_theta = A(4:6, 7:9)      % a pitch error pushes along +x, a roll error along -y
B_v_T     = B(4:6, 1)'       % extra thrust only changes vertical acceleration (1/m)
B_w_tau   = B(10:12, 2:4)    % torque -> angular acceleration (J^-1)
\end{toycode}

\begin{real}
Our Python node does not compute $A$ and $B$; they are used \emph{inside the estimator} (block 4), which in our
SITL setup is ArduPilot's EKF3. They are also the reason Model~2 can use a linear model at all: the table above shows
that near hover the drone really is almost linear. And they are what Model~2 learns from data
(\cref{sec:m2b2}): DMDc recovers these same numbers without knowing $m$ or $J$.
\end{real}

% ---------------------------------------------------------------------
\subsection{Block 3 --- How a small error grows (local error dynamics)}
\label{sec:m1b3}
% ---------------------------------------------------------------------

\begin{idea}
Suppose our belief about the drone is slightly wrong. How does that mistake change over time if we only trust the
physics? Block 3 answers this with the linear model of block 2. This is exactly what the estimator needs to know to
decide how much it can still trust its own prediction.
\end{idea}

\begin{eqbox}{Equations (7)--(9): linear dynamics of the error}
\begin{align}
  \delta\vx &= \big[\vp-\hp;\ \vv-\hv;\ \Log(\hR^\top R);\ \vw-\hw\big] \tag{7}\label{eq:7}\\
  \delta\dot{\vx} &= A\,\delta\vx + B\,\delta\vu \tag{8}\label{eq:8}\\
  \delta\vx_{k+1} &= F\,\delta\vx_k + G\,\delta\vu_k,\qquad F=e^{A\Delta t}\approx I+A\Delta t,\quad G\approx B\Delta t \tag{9}\label{eq:9}
\end{align}
\end{eqbox}

\subsubsection*{Derivation}
The truth follows $\dot\vx=\vf(\vx,\vu)$ and the estimate follows the same model, $\dot\hx=\vf(\hx,\hat\vu)$.
Subtract the two and use (5):
\[
\delta\dot\vx=\vf(\vx,\vu)-\vf(\hx,\hat\vu)\approx\big[\vf(\hx,\hat\vu)+A\delta\vx+B\delta\vu\big]-\vf(\hx,\hat\vu)=A\delta\vx+B\delta\vu .
\]
The big nonlinear parts cancel; only the small linear error is left. That is (8).

For a computer we need one time step $\Delta t$. The exact solution of $\delta\dot\vx=A\delta\vx$ over one step is
$\delta\vx_{k+1}=e^{A\Delta t}\delta\vx_k$, where the matrix exponential is the series
$e^{A\Delta t}=I+A\Delta t+\frac12(A\Delta t)^2+\dots$ (Taylor again). Keeping the first two terms gives
$F\approx I+A\Delta t$, which is the same as one Euler step.

\textbf{Why a drone drifts away.} At hover every eigenvalue of $A$ is 0. Look at the chain inside $A$:
angular-velocity error $\to$ attitude error $\to$ velocity error $\to$ position error. Each arrow is an
integration. A constant roll-rate error $\delta\omega_x$ therefore becomes
$\delta\theta_x=\delta\omega_x t$, then $\delta v_y=-g_0\delta\omega_x t^2/2$, then
\[
\delta y(t)=-\frac{g_0\,\delta\omega_x\,t^3}{6}.
\]
\begin{toy}[a tiny gyro error becomes a big position error]
$\delta\omega_x=0.01$\,rad/s (about half a degree per second, a typical gyro bias). After 2\,s:
$\delta y=-9.81\times0.01\times8/6=-0.131$\,m. After 5\,s: $-2.04$\,m. After 10\,s: $-16.4$\,m. This $t^3$ growth is why an
IMU on its own is useless for position after a few seconds, and why the camera is needed.

\textbf{Linear prediction versus truth.} The toy model starts the drone with an initial error (10\,cm, 0.2\,rad/s of
yaw rate, and a roll error) and compares the linear prediction $e^{At}\delta\vx_0$ with the true nonlinear motion over
2\,s. With a 2$^\circ$ roll error they differ by only 0.02\%; with a 30$^\circ$ roll error the linear model is off
by 4.7\% (\cref{fig:errgrow}). The linear model is \emph{local}: perfect for small errors, worse for large ones.

\textbf{One-step check.} With $\Delta t=0.01$\,s, the biggest entry of $F-e^{A\Delta t}$ is $4.9\times10^{-4}$,
the size of the dropped $\tfrac12(A\Delta t)^2$ term. So $F=I+A\Delta t$ is good enough.
\end{toy}

\begin{figure}[H]
  \centering
  \includegraphics[width=0.85\linewidth]{M1_1_Dynamics_Linearization_2.png}
  \caption{Sideways error $\delta y$ after an initial roll error: nonlinear truth (solid) and linear model (dashed).}
  \label{fig:errgrow}
\end{figure}

\begin{toycode}{M1\_1\_Dynamics\_Linearization}{Block 3}
dt = 0.01;                          % EKF / IMU sample time [s]
F  = eye(12) + A*dt;
max_F_error = max(abs(F - expm(A*dt)), [], 'all')   % size of the dropped (A dt)^2/2
eig_A = eig(A)'                                       % all zero at hover
\end{toycode}

\begin{real}
The estimator never tries to describe the whole flight with a linear model. It only describes the \emph{small
difference} between what it believes and what is true, and that difference stays small because the camera and the
IMU keep correcting it. That is why the filter of block~4 is called an \emph{error-state} filter, and why $F$ is all
it needs.
\end{real}

% ---------------------------------------------------------------------
\subsection{Block 4 --- Where am I? The error-state Kalman filter (our VIO)}
\label{sec:m1b4}
% ---------------------------------------------------------------------

\begin{idea}
We have two imperfect sources of information. The \emph{physics} (blocks 1--3) predicts where the drone should be,
but small mistakes grow quickly. The \emph{sensors} tell us where it is, but every reading is noisy, and the camera
only gives a new reading 25 times a second. The Kalman filter combines them in the best possible way. It repeats
two steps for ever:
\begin{enumerate}
  \item \textbf{Predict.} ``I know how I pushed the drone, so physics tells me where it should be now.'' The guess
        moves forward, and because physics is never perfect, the uncertainty \emph{grows}.
  \item \textbf{Correct.} ``A sensor reading arrived. How surprised am I?'' The guess moves towards the reading. How
        far depends on whom we trust more. The uncertainty \emph{shrinks}.
\end{enumerate}
The filter remembers two things: the \textbf{estimate} $\hx_k$ (its best guess) and the \textbf{covariance} $P_k$
(a $12\times12$ matrix that says how unsure it is in each of the 12 error directions).
\end{idea}

\begin{eqbox}{Equations (10)--(14): error-state extended Kalman filter}
\textbf{Prediction (every IMU sample, 100\,Hz)}
\begin{align}
  \hx_{k+1} &= \vf_d(\hx_k,\vu_k) \tag{10}\label{eq:10}\\
  P_{k+1}   &= F P_k F^\top + Q \tag{11}\label{eq:11}
\end{align}
\textbf{Correction (whenever a measurement $\vz$ arrives)}
\begin{align}
  K   &= P H^\top\big(HPH^\top+R_n\big)^{-1} \tag{12}\label{eq:12}\\
  \delta\hat{\vx} &= K\big(\vz-h(\hx)\big),\quad
  \hp\leftarrow\hp+\delta\hat{\vp},\ \hv\leftarrow\hv+\delta\hat{\vv},\ \hR\leftarrow\hR\,\Exp(\delta\hat{\vth}),\ \hw\leftarrow\hw+\delta\hat{\vw} \tag{13}\label{eq:13}\\
  P   &\leftarrow (I-KH)\,P \tag{14}\label{eq:14}
\end{align}
$\vf_d$ = one simulation step of (1)--(4); $Q$ = how much the physics may be wrong per step (wind, model errors);
$\vz$ = measurement; $h(\hx)$ = what we expect to measure if our estimate were right; $H$ = Jacobian of $h$;
$R_n$ = covariance of the sensor noise; $K$ = Kalman gain.
\end{eqbox}

\subsubsection*{First understand it with one number}
Before matrices, take a single quantity, the altitude $z$, with a guess $\hat z$ and variance $P$.

\textbf{Predict.} If we command a climb of $c$\,m/s for $\Delta t$ seconds, the guess becomes $\hat z+c\Delta t$. The old
uncertainty is still there, and the physics adds a bit more ($Q$): $P^-=P+Q$.

\textbf{Correct.} A sensor reads $z_m$ with noise variance $R_n$. The surprise (the \textbf{innovation}) is
$r=z_m-\hat z^-$. We move a fraction $K$ of the surprise: $\hat z^+=\hat z^-+K\,r$. Which fraction is best? The new
error is $(1-K)\times$(old error) $-K\times$(sensor noise), whose variance is
\[
P^+=(1-K)^2P^-+K^2R_n .
\]
This is a bowl in $K$ (\cref{sec:minquad}). Set the slope to zero:
$-2(1-K)P^-+2KR_n=0\Rightarrow K=\dfrac{P^-}{P^-+R_n}$. Putting it back gives $P^+=(1-K)P^-$.

So $K$ is simply ``my uncertainty divided by (my uncertainty + the sensor's uncertainty)''. If I am very unsure
($P^-$ large), $K\to1$ and I believe the sensor. If the sensor is very noisy ($R_n$ large), $K\to0$ and I ignore it.

\begin{toy}[one predict and one correct, altitude only]
Prior: $\hat z_k=10$\,m with $P_k=3.5$\,m$^2$. We command a climb of 0.5\,m/s for 1\,s, $F=1$, $Q=0.5$:
\[
\hat z^-=10+0.5=10.5\ \text{m},\qquad P^-=1\cdot3.5\cdot1+0.5=4.0\ \text{m}^2 .
\]
A sensor reads $z=12.5$\,m with $R_n=1$\,m$^2$ and $H=1$:
\[
K=\frac{4.0}{4.0+1}=0.8,\qquad \hat z^+=10.5+0.8\,(12.5-10.5)=12.1\ \text{m},\qquad P^+=(1-0.8)\cdot4.0=0.8\ \text{m}^2 .
\]
The sensor is four times more trustworthy than the prediction ($4.0$ versus $1$), so the estimate moves 80\% of the
way towards it, and the variance drops from 4.0 to 0.8. Both sources together are better than either alone.
\end{toy}

\begin{toycode}{M1\_2\_Error\_State\_EKF}{Toy example - altitude only}
z_hat = 10;  Pk = 3.5;  F = 1;  Q = 0.5;  climb = 0.5;  dt = 1;
z_hat = z_hat + climb*dt          % predict the guess
Pk    = F*Pk*F' + Q               % uncertainty grows
z = 12.5;  H = 1;  Rn = 1;
K     = Pk*H'/(H*Pk*H' + Rn)      % how much to trust the sensor
z_hat = z_hat + K*(z - H*z_hat)   % corrected guess
Pk    = (1 - K*H)*Pk              % uncertainty shrinks
\end{toycode}

\subsubsection*{The same thing with matrices}
Now the error is the 12-vector $\delta\vx$ and the uncertainty is the matrix $P$.

\textbf{(10) Predict the guess.} The best guess of the next state is to run the full nonlinear physics one step
forward from the current guess with the input we applied. The error we do not know averages to zero, so it does not
move the guess.

\textbf{(11) Predict the uncertainty.} From block 3, the error moves as $\delta\vx_{k+1}=F\delta\vx_k+\bm w_k$, where
$\bm w_k$ is the new random error added by the wind and by model mistakes, with covariance $Q$. Using the rule
$P'=FPF^\top$ of \cref{sec:maths} and the fact that $\bm w_k$ is independent of the old error:
\[
P_{k+1}=\mathbb E[\delta\vx_{k+1}\delta\vx_{k+1}^\top]=F\,P_k\,F^\top+Q .
\]
Uncertainty is stretched by the physics ($FPF^\top$) and inflated by the noise ($+Q$).

\textbf{(12) The Kalman gain.} A sensor does not see the whole state. A camera sees position and attitude; a gyro
sees angular velocity. $H$ picks out what the sensor sees: the innovation is $\vr=\vz-h(\hx)\approx H\delta\vx+\bm v$
(``what I see minus what I expected = my error seen through the sensor + sensor noise''). Correcting by
$\delta\hat\vx=K\vr$ leaves the error $(I-KH)\delta\vx-K\bm v$ with covariance
\[
P^+=(I-KH)P(I-KH)^\top+KR_nK^\top .
\]
Exactly as in the one-number case, we choose $K$ to make the total uncertainty (the sum of the diagonal, the trace)
as small as possible. Setting the derivative with respect to $K$ to zero:
\[
-2PH^\top+2K\big(HPH^\top+R_n\big)=0\quad\Longrightarrow\quad K=PH^\top\big(HPH^\top+R_n\big)^{-1}.
\]
Compare with $K=P^-/(P^-+R_n)$: it is the same formula, with matrices.

\textbf{(13) Correct the guess.} $\delta\hat\vx=K\vr$ is our best estimate of the error, so we add it to the guess.
Position, velocity and angular velocity are added normally. The attitude is \emph{rotated}: $\hR\leftarrow\hR\Exp(\delta\hat\vth)$,
so $\hR$ stays a true rotation matrix. This is the ``error-state'' trick: the filter works with a small 3-number
attitude error and never adds matrices that should be rotations.

\textbf{(14) Shrink the uncertainty.} Putting the best $K$ into the expression for $P^+$ simplifies it to
$P^+=(I-KH)P$.

\subsubsection*{Our sensors}
\begin{center}\small
\begin{tabular}{@{}llll@{}}
\toprule
Sensor (rate) & innovation $\vr$ & $H$ & noise $\sigma$\\
\midrule
Gyroscope (100\,Hz) & $\tilde{\vw}-\hw$ & $[\,0\ \ 0\ \ 0\ \ I\,]$ & 0.01\,rad/s\\
Camera / VIO (25\,Hz) & $\begin{bmatrix}\tilde{\vp}-\hp\\ \Log(\hR^\top\tilde R)\end{bmatrix}$ &
  $\begin{bmatrix}I&0&0&0\\0&0&I&0\end{bmatrix}$ & 5\,cm, 0.02\,rad\\
\bottomrule
\end{tabular}
\end{center}
Each block of $H$ is $3\times3$, in the order $[\delta\vp,\delta\vv,\delta\vth,\delta\vw]$. The \textbf{velocity is never
measured}. The filter works it out from how the measured position changes, through the physics in $F$.

\begin{toy}[the full 12-state filter on the drone (\code{M1\_2})]
The drone flies a climbing helix in gusty wind for 20\,s. The filter starts with a wrong guess (30\,cm, 0.2\,m/s and
about 7$^\circ$ off) and sees only the commands and the noisy sensors. After the first 2\,s:
\begin{center}\small
\begin{tabular}{lrl}
\toprule
quantity & estimation error (RMS) & for comparison\\
\midrule
position & 2.55\,cm & the camera alone: 8.55\,cm\\
velocity (never measured!) & 3.85\,cm/s & --\\
attitude & 0.20$^\circ$ & camera noise: 1.15$^\circ$\\
\bottomrule
\end{tabular}
\end{center}
The fused estimate is more than three times better than the camera alone. In 99\% of the samples all 12 errors lie
inside the $\pm3\sigma$ band that the filter itself predicts from $P_k$: the filter knows how unsure it is.
\end{toy}

\begin{figure}[H]
  \centering
  \includegraphics[width=0.8\linewidth]{M1_2_Error_State_EKF_1.png}\\[4pt]
  \includegraphics[width=0.9\linewidth]{M1_2_Error_State_EKF_2.png}
  \caption{Top: noisy camera positions (grey dots), truth (black) and EKF estimate (red, dashed). Bottom: estimation
  errors (blue) stay inside the $\pm3\sigma$ band computed from $P_k$.}
\end{figure}

\begin{toycode}{M1\_2\_Error\_State\_EKF}{Run the filter}
for k = 1:N
    [~, ~, ~, wh] = DroneLib.unpack(xh);                       % gyro correction
    [xh, Pk] = DroneLib.ekf_update(xh, Pk, Zg(:,k) - wh, Hg, Rg);
    if mod(k-1, 4) == 0                                % camera (25 Hz)
        [ph, ~, Rh] = DroneLib.unpack(xh);
        r = [Zp(:,k) - ph; DroneLib.logSO3(Rh'*DroneLib.expSO3(Zth(:,k)))];
        [xh, Pk] = DroneLib.ekf_update(xh, Pk, r, Hc, Rc);
    end
    ...
    [xh, Pk] = DroneLib.ekf_predict(xh, Pk, U(:,k), dt, P, Q);  % prediction
end
\end{toycode}
\begin{lstlisting}[style=mat, title={\small\sffamily toy\_model3/DroneLib.m --- ekf\_predict and ekf\_update, equations (9)--(14)}]
A  = DroneLib.jacobians(xh, u, P);
F  = eye(12) + A*dt;                                 % (9)
xh = DroneLib.rk4(xh, u, dt, P);                     % (10)
Pk = F*Pk*F' + Q;                                    % (11)

K  = Pk*H' / (H*Pk*H' + Rn);                         % (12)
xh = DroneLib.boxplus(xh, K*r);                      % (13)  R <- R*Exp(dtheta)
Pk = (eye(12) - K*H)*Pk;                             % (14)
\end{lstlisting}

\begin{real}[this is the ``VIO'' of our title]
A monocular camera alone cannot tell scale (a small room close up looks like a big room far away) and an IMU alone
drifts within seconds ($t^3$, block 3). The Kalman filter is the referee that fuses them. Between camera frames it
predicts with the physics and the gyro; when a frame arrives, the visual-odometry pose corrects the prediction, and
$K$ decides how much to trust each source.

In our SITL setup this filter runs \emph{inside ArduPilot} (its EKF3 has exactly this predict / correct structure,
with more states such as sensor biases). Its output arrives on \code{/ap/v1/pose/filtered} and
\code{/ap/v1/twist/filtered}. Our node stores it, and before every control step checks that the estimate is complete,
recent and finite:
\pyfile{447}{480}{\node}{drones\_controller/controller\_node.py --- state\_is\_valid() (first part)}
If the estimate is older than \code{state\_timeout} = 0.5\,s, the controller does nothing at all. This is a safety rule:
never fly on stale information.
\end{real}

% ---------------------------------------------------------------------
\subsection{Block 5 --- Position control (outer loop)}
\label{sec:m1b5}
% ---------------------------------------------------------------------

\begin{idea}
Compare where we are with where we should be, and ask for an acceleration that pulls us there: the further away,
the harder we pull (a spring), and the faster we are already moving towards the target, the less we pull (a damper).
\end{idea}

\begin{eqbox}{Equations (15)--(16): tracking errors and desired acceleration}
\begin{align}
  \ve_p=\vp_d-\hp,\qquad \ve_v&=\vv_d-\hv \tag{15}\label{eq:15}\\
  \va_d&=\ddot{\vp}_d+K_p\,\ve_p+K_v\,\ve_v \tag{16}\label{eq:16}
\end{align}
$\vp_d,\vv_d,\ddot\vp_d$: desired position, velocity, acceleration. $K_p,K_v$: diagonal gain matrices.
\end{eqbox}

\subsubsection*{Why this works: a spring and a damper}
Suppose the inner loops deliver exactly the acceleration we ask for, $\ddot{\vp}=\va_d$. Then the position error
$\ve_p=\vp_d-\vp$ changes as
\[
\ddot{\ve}_p=\ddot{\vp}_d-\ddot{\vp}=\ddot{\vp}_d-\va_d=-K_p\ve_p-K_v\dot{\ve}_p
\quad\Longrightarrow\quad \ddot{\ve}_p+K_v\dot{\ve}_p+K_p\ve_p=\bm0 .
\]
This is the equation of a mass on a spring (stiffness $k_p$) with a damper (friction $k_v$), one per axis. Its
solutions behave like $e^{st}$ with $s^2+k_vs+k_p=0$. If $k_p>0$ and $k_v>0$, both roots have negative real part, so the
error always dies out. Comparing with the standard form $s^2+2\zeta\omega_ns+\omega_n^2$:
\[
\omega_n=\sqrt{k_p}\ \text{(how fast)},\qquad \zeta=\frac{k_v}{2\sqrt{k_p}}\ \text{(how damped)},\qquad t_{\text{settle}}\approx\frac{4}{\zeta\omega_n}.
\]
\begin{center}\small
\begin{tabular}{@{}lrrrrr@{}}
\toprule
axis & $k_p$ & $k_v$ & $\omega_n$ [rad/s] & $\zeta$ & $t_{\text{settle}}$ [s]\\
\midrule
$x,\ y$ & 1.5 & 1.2 & 1.225 & 0.490 & 6.7\\
$z$ & 2.0 & 1.5 & 1.414 & 0.530 & 5.3\\
\bottomrule
\end{tabular}
\end{center}
Our \code{controller.yaml} gains are gentle: $\zeta\approx0.5$ means a small overshoot, and the error settles in
about 6\,s. That suits a camera drone that should move smoothly. The term $\ddot{\vp}_d$ (``feed-forward'') supplies the
acceleration a perfect pilot would need to follow a moving target; the feedback terms only fix what is left over.

\begin{toy}[50 cm too low and already climbing]
Setpoint $\vp_d=[0,0,2]$\,m (the \code{controller.yaml} hover point), $\vv_d=\ddot{\vp}_d=\bm0$. The estimator says
$\hp=[0.1,\,-0.2,\,1.5]$\,m and $\hv=[0,0,0.3]$\,m/s.
\[
\ve_p=\begin{bmatrix}0-0.1\\0-(-0.2)\\2-1.5\end{bmatrix}=\begin{bmatrix}-0.1\\0.2\\0.5\end{bmatrix},\quad
\ve_v=\begin{bmatrix}0\\0\\-0.3\end{bmatrix},\quad
\va_d=\begin{bmatrix}1.5(-0.1)\\1.5(0.2)\\2(0.5)+1.5(-0.3)\end{bmatrix}=\begin{bmatrix}-0.15\\0.30\\0.55\end{bmatrix}\text{ m/s}^2 .
\]
It is 50\,cm low, so it wants to accelerate up. Because it is already climbing at 0.3\,m/s, it asks for less: 0.55
instead of 1.0. That is the damper at work, preventing overshoot.
\end{toy}

\begin{toycode}{M1\_3\_SO3\_Control\_Closed\_Loop}{Block 5}
p_hat = [0.1; -0.2; 1.5];   v_hat = [0; 0; 0.3];
p_d   = [0; 0; 2];          v_d = [0; 0; 0];     a_ff = [0; 0; 0];
e_p = p_d - p_hat
e_v = v_d - v_hat
a_d = a_ff + P.Kp*e_p + P.Kv*e_v
\end{toycode}

\begin{real}
Equation (15) is \code{compute\_position\_errors()} and equation (16) is the first part of
\code{compute\_desired\_force()} in our controller (\code{a\_cmd} is $\va_d$):
\pyfile{32}{40}{\so}{drones\_controller/so3\_controller.py --- equation (15)}
\pyfile{64}{68}{\so}{drones\_controller/so3\_controller.py --- equation (16)}
The desired state comes from \code{WaypointGenerator}: a fixed point $[0,0,2]$ with $\vv_d=\va_d=\bm0$, read from the YAML.
\end{real}

% ---------------------------------------------------------------------
\subsection{Block 6 --- Force and attitude: which way to lean}
\label{sec:m1b6}
% ---------------------------------------------------------------------

\begin{idea}
A quadrotor can only push along its own up-axis. To get the acceleration $\va_d$ it must (1) work out the total force
needed, including holding up its own weight, and (2) tilt so that its up-axis points along that force. Like a person
pushing a heavy cart: to push forward, you lean forward.
\end{idea}

\begin{eqbox}{Equations (17)--(19): desired force, thrust direction and attitude}
\begin{align}
  \vF_d &= m\,(\va_d-\vg) \tag{17}\label{eq:17}\\
  \vb_{3d} &= \frac{\vF_d}{\|\vF_d\|} \tag{18}\label{eq:18}\\
  \vb_{1c}=[\cos\psi_d,\ \sin\psi_d,\ 0]^\top,\quad
  \vb_{2d}&=\frac{\vb_{3d}\times\vb_{1c}}{\|\vb_{3d}\times\vb_{1c}\|},\quad
  \vb_{1d}=\vb_{2d}\times\vb_{3d},\quad
  R_d=\big[\vb_{1d}\ \ \vb_{2d}\ \ \vb_{3d}\big] \tag{19}\label{eq:19}
\end{align}
\end{eqbox}

\subsubsection*{Derivation}
\textbf{(17)} We want $\dot\vv=\va_d$. Equation (2) says $\dot\vv=\vg+\frac{T}{m}R\ve_3$, so the thrust vector must be
$T\,R\ve_3=m(\va_d-\vg)$. Call this $\vF_d$. Because $\vg$ points down, $-\vg$ points up: even for $\va_d=0$ the drone
must push up with $mg_0=14.7$\,N.

\textbf{(18)} $T R\ve_3=T\vb_3$ is a vector of length $T$ along the drone's up-axis. For it to equal $\vF_d$, the up-axis
must point along $\vF_d$, so $\vb_{3d}=\vF_d/\|\vF_d\|$, and the thrust must be $T=\|\vF_d\|$.
This is the key fact of quadrotor flight: \emph{position is controlled through the attitude}.

\textbf{(19)} $\vb_{3d}$ fixes the tilt (two of the three rotation directions). The third, yaw (where the nose points),
is our free choice $\psi_d$. We build the remaining two axes so that they are unit length and at right angles:
\begin{itemize}
  \item $\vb_{1c}$ is the heading we would like, flat on the ground.
  \item $\vb_{2d}=\vb_{3d}\times\vb_{1c}$ (normalised) is at right angles to both the up-axis and the heading: the drone's ``left''.
  \item $\vb_{1d}=\vb_{2d}\times\vb_{3d}$ is at right angles to both: the nose, as close to the wanted heading as the tilt allows.
\end{itemize}
The three columns are orthonormal and right-handed, so $R_d^\top R_d=I$ and $\det R_d=1$: $R_d$ is a true rotation.
The construction only fails if $\vb_{3d}$ is parallel to $\vb_{1c}$ (a 90$^\circ$ tilt); the code then nudges the yaw.

\begin{toy}[continuing the example of block 5]
$\va_d=[-0.15,\,0.30,\,0.55]$, $m=1.5$, $\vg=[0,0,-9.81]$, $\psi_d=0$:
\[
\vF_d=1.5\,[-0.15,\ 0.30,\ 0.55+9.81]=[-0.225,\ 0.450,\ 15.540]\ \text{N},\qquad \|\vF_d\|=15.548\ \text{N},
\]
\[
\vb_{3d}=[-0.0145,\ 0.0289,\ 0.9995],\qquad
R_d=\begin{bmatrix}0.9999&0.0000&-0.0145\\0.0004&0.9996&0.0289\\0.0145&-0.0289&0.9995\end{bmatrix}.
\]
Tilt $=\arccos(0.9995)=1.85^\circ$: lean very slightly towards $-x$ and $+y$ (where the target is) and push 15.5\,N,
more than the 14.7\,N weight because the drone must also climb. Check: $\|R_d^\top R_d-I\|=1.1\times10^{-16}$.
\end{toy}

\begin{toycode}{M1\_3\_SO3\_Control\_Closed\_Loop}{Block 6}
F_d = P.m*(a_d - P.g)                 % N
b3d = F_d/norm(F_d);
b1c = [cos(psi_d); sin(psi_d); 0];
b2d = cross(b3d, b1c)/norm(cross(b3d, b1c));
b1d = cross(b2d, b3d);
R_d = [b1d b2d b3d]
\end{toycode}

\begin{real}
Equation (17) is the end of \code{compute\_desired\_force()}. The code writes $\vF_d=m(\va_{\text{cmd}}+g\ve_3)$, which is
the same thing because $-\vg=g\ve_3$ in ENU:
\pyfile{70}{74}{\so}{drones\_controller/so3\_controller.py --- equation (17)}
Equations (18)--(19) are \code{compute\_desired\_attitude()}, including the safety nudge of the yaw:
\pyfile{82}{126}{\so}{drones\_controller/so3\_controller.py --- equations (18)--(19)}
The node converts $R_d$ to a quaternion and publishes it on the topic\\
\code{/drones\_controller/desired\_attitude}, so we can watch in RViz which way the controller wants to lean.
\end{real}

% ---------------------------------------------------------------------
\subsection{Block 7 --- Attitude control on SO(3) (inner loop)}
\label{sec:m1b7}
% ---------------------------------------------------------------------

\begin{idea}
``I should be leaning like $R_d$, but I am leaning like $\hR$. Which way do I twist, and how hard?'' The inner loop
computes an attitude error directly from the two rotation matrices (no angles, so no gimbal lock), a spin error, and
a torque that removes both, like a spring and a damper for rotation. It runs faster than the position loop, because
the drone must tilt before it can move.
\end{idea}

\begin{eqbox}{Equations (20)--(23): geometric attitude controller}
\begin{align}
  \ve_R &= \tfrac12\big(R_d^\top\hR-\hR^\top R_d\big)^\vee \tag{20}\label{eq:20}\\
  \ve_\omega &= \hw-\hR^\top R_d\,\vw_d \tag{21}\label{eq:21}\\
  \vtau &= -K_R\ve_R-K_\omega\ve_\omega+\hw\times J\hw
          -J\big(\skw{\hw}\hR^\top R_d\vw_d-\hR^\top R_d\dot{\vw}_d\big) \tag{22}\label{eq:22}\\
  T &= \vF_d\cdot\hR\,\ve_3 \tag{23}\label{eq:23}
\end{align}
\end{eqbox}

\subsubsection*{(20) The attitude error}
We need a number that says how far apart two attitudes are. Use
\[
\Psi=\tfrac12\tr\big(I-R_d^\top\hR\big)=1-\cos\varphi ,
\]
where $\varphi$ is the angle of the rotation that takes $R_d$ to $\hR$ (because the trace of a rotation by $\varphi$ is
$1+2\cos\varphi$). $\Psi=0$ when the attitudes agree and $\Psi=2$ when they are upside down to each other.
Now turn the drone a tiny bit in the direction $\bm\eta$ and see how $\Psi$ changes. A short calculation gives
\[
\text{change of }\Psi=\ve_R\cdot\bm\eta,\qquad \ve_R=\tfrac12\big(R_d^\top\hR-\hR^\top R_d\big)^\vee .
\]
So $\ve_R$ is the \emph{slope} (gradient) of the error: it points in the direction of turning that increases the
error fastest. The controller twists the opposite way, $-K_R\ve_R$, ``downhill''. For a single-axis error of angle
$\varphi$ about the axis $\bm n$, $\ve_R=\sin\varphi\,\bm n\approx\varphi\,\bm n$ for small angles.

\subsubsection*{(21) The angular-velocity error}
The relative rotation $R_e=R_d^\top\hR$ changes as $\dot R_e=R_e\skw{\ve_\omega}$ with $\ve_\omega=\hw-\hR^\top R_d\vw_d$
(differentiate and use (3) for both). So $\ve_\omega$ is how fast the drone turns \emph{relative to how it should be
turning}, in body axes. $\hR^\top R_d\vw_d$ is simply the desired spin written in the drone's own axes.
Note the order: \emph{actual minus desired}.

\subsubsection*{(22) The torque, and why it is guaranteed to work}
Differentiate $\ve_\omega$ and replace $J\dot{\hw}$ using (4):
\[
J\dot{\ve}_\omega=\vtau-\hw\times J\hw+J\big(\skw{\hw}\hR^\top R_d\vw_d-\hR^\top R_d\dot{\vw}_d\big).
\]
Now insert the torque (22). The gyroscopic term and the bracket cancel exactly, and what remains is
\[
J\dot{\ve}_\omega=-K_R\ve_R-K_\omega\ve_\omega ,
\]
a rotational spring ($K_R$) and damper ($K_\omega$). To prove it always settles, take (with $K_R=k_RI$, $K_\omega=k_\omega I$)
the ``energy''
\[
V=\underbrace{\tfrac12\ve_\omega^\top J\ve_\omega}_{\text{spin energy}}+\underbrace{k_R\Psi}_{\text{spring energy}} .
\]
Its rate of change is $\dot V=\ve_\omega^\top(-k_R\ve_R-k_\omega\ve_\omega)+k_R\,\ve_R\cdot\ve_\omega=-k_\omega\|\ve_\omega\|^2\le0$.
The energy can only go down, like a pendulum with friction, so the drone settles at $R_d$ from \emph{any} starting
attitude except exactly upside down. This kind of proof is called a \textbf{Lyapunov} argument.

\subsubsection*{(23) The thrust}
Thrust can only act along the \emph{current} up-axis $\hR\ve_3$. The useful part of $\vF_d$ along that axis is the dot
product $T=\vF_d\cdot\hR\ve_3$. When the attitude is right, $T=\|\vF_d\|$. While the drone is still tilted the wrong way,
$T$ is automatically smaller, so it does not push hard in a wrong direction.

\begin{toy}[recovering from a 10 degree roll]
The drone is rolled $10^\circ$ and still rolling further at 0.2\,rad/s; it should be level ($R_d=I$, $\vw_d=\dot\vw_d=0$),
$\vF_d=[0,0,14.715]$:
\[
\ve_R=\tfrac12(\hR-\hR^\top)^\vee=[\sin10^\circ,\,0,\,0]=[0.1736,\,0,\,0],\qquad \ve_\omega=[0.2,\,0,\,0],
\]
\[
\tau_x=-4(0.1736)-0.8(0.2)+0=-0.8546\ \text{N\,m},\qquad T=14.715\cos10^\circ=14.491\ \text{N}.
\]
A torque about $-x$ rolls it back, and the thrust is reduced to the part that is still vertical.

\textbf{Recovery from 60$^\circ$.} The toy model starts the drone tilted 60$^\circ$ and lets block 7 bring it back. After 3\,s
the error is $2\times10^{-6}$ degrees, and the energy $V$ goes down all the time (\cref{fig:att}).
\end{toy}

\begin{figure}[H]
  \centering
  \includegraphics[width=0.85\linewidth]{M1_3_SO3_Control_Closed_Loop_1.png}
  \caption{Block 7 alone: from 60$^\circ$ to level. Right: the energy $V$ only goes down, as the proof promised.}
  \label{fig:att}
\end{figure}

\begin{toycode}{M1\_3\_SO3\_Control\_Closed\_Loop}{Block 7}
e_R = 0.5*DroneLib.vee(Rd'*R_hat - R_hat'*Rd)
e_w = w_hat - R_hat'*Rd*w_d
tau = -P.KR*e_R - P.Kw*e_w + cross(w_hat, P.J*w_hat) ...
      - P.J*(DroneLib.hat(w_hat)*R_hat'*Rd*w_d - R_hat'*Rd*wdot_d)
T = dot(Fd, R_hat*P.e3)
\end{toycode}

\begin{real}
Equations (20), (21) and (22) are three functions of our controller, written exactly as above:
\pyfile{134}{142}{\so}{drones\_controller/so3\_controller.py --- equation (20)}
\pyfile{152}{158}{\so}{drones\_controller/so3\_controller.py --- equation (21)}
\pyfile{181}{210}{\so}{drones\_controller/so3\_controller.py --- equation (22)}
Equation (23) is not in the Python code, because our node does not send thrust to ArduPilot (see
\cref{sec:realctrl}). The toy model computes it.
\end{real}

% ---------------------------------------------------------------------
\subsection{Block 8 --- Four motor speeds (control allocation)}
\label{sec:m1b8}
% ---------------------------------------------------------------------

\begin{idea}
The drone does not have a ``torque knob''. It has four motors. Spinning the right-side motors faster rolls it left;
spinning the front ones faster pitches it back; spinning the clockwise pair faster yaws it. Block 8 solves the
reverse problem: which four rotor thrusts give exactly the wanted $T$ and $\vtau$?
\end{idea}

\begin{eqbox}{Equations (24)--(25): control allocation and motor speeds}
\begin{align}
  \begin{bmatrix}T\\ \vtau\end{bmatrix}&=\Gamma\,\vf,\qquad
  \Gamma=\begin{bmatrix}1&1&1&1\\ y_1&y_2&y_3&y_4\\ -x_1&-x_2&-x_3&-x_4\\ \kappa_1c_m&\kappa_2c_m&\kappa_3c_m&\kappa_4c_m\end{bmatrix} \tag{24}\label{eq:24}\\
  \Omega_i&=\sqrt{\frac{\big(\Gamma^{-1}\vu\big)_i}{k_f}}\quad\text{clipped to }[\Omega_{\min},\Omega_{\max}] \tag{25}\label{eq:25}
\end{align}
$f_i$: thrust of rotor $i$; $(x_i,y_i)$: its position on the frame; $\kappa_i=\pm1$: its spin direction; $c_m$: yaw torque per newton;
$k_f$: thrust constant ($f_i=k_f\Omega_i^2$).
\end{eqbox}

\subsubsection*{Derivation}
Each rotor $i$ pushes up with $f_i$ at the point $\bm r_i=[x_i,y_i,0]$ of the frame.
\begin{itemize}
  \item Total thrust: the pushes add up, $T=f_1+f_2+f_3+f_4$ (row 1 of $\Gamma$).
  \item Roll and pitch torque: a force at a distance creates a torque $\bm r_i\times f_i\ve_3=f_i[y_i,\,-x_i,\,0]$
        (rows 2 and 3). A rotor on the right ($y_i<0$) creates negative roll torque.
  \item Yaw torque: air resists each spinning propeller, so each motor twists the body the opposite way with a torque
        proportional to its thrust, $\kappa_ic_mf_i$ (row 4).
\end{itemize}
A propeller's thrust grows with the square of its speed: $f_i=k_f\Omega_i^2$. Inverting both relations gives (25).
Motors cannot spin backwards or faster than their maximum, so the speeds are clipped.

For the Iris (ArduPilot ``Quad-X'': 1 front-right CCW, 2 back-left CCW, 3 front-left CW, 4 back-right CW):
\[
\Gamma=\begin{bmatrix}1&1&1&1\\-0.22&0.22&0.22&-0.22\\-0.13&0.13&-0.13&0.13\\-0.016&-0.016&0.016&0.016\end{bmatrix}.
\]

\begin{toy}[hover, and the roll correction of block 7]
\textbf{Hover} $\vu=[14.715,0,0,0]$: by symmetry every rotor gives $f_i=14.715/4=3.679$\,N, so
$\Omega_i=\sqrt{3.679/8.549\times10^{-6}}=656$\,rad/s (about 6260\,rpm).

\textbf{Roll correction} $\vu=[14.491,\,-0.8546,\,0,\,0]$: by symmetry $f_1=f_4=a$ and $f_2=f_3=b$. Row 1:
$2a+2b=14.491$. Row 2: $-0.22a+0.22b+0.22b-0.22a=0.44(b-a)=-0.8546$. Solving: $a=4.594$\,N, $b=2.652$\,N, so
$\vOmega=[733,\ 557,\ 557,\ 733]$\,rad/s. The right-side rotors (1 and 4) speed up, the left ones slow down, and the
drone rolls back to level.
\end{toy}

\begin{toycode}{M1\_3\_SO3\_Control\_Closed\_Loop}{Block 8 (DroneLib.mixer)}
fr = P.Gamma \ u;                                          % rotor thrusts
fr = min(max(fr, P.kf*P.Omega_min^2), P.kf*P.Omega_max^2); % motor limits
Omega = sqrt(fr/P.kf);                                     % motor speeds
\end{toycode}

\begin{real}
This block lives inside ArduPilot (its ``motor mixer''); our node never talks to the motors directly. The matrix
$\Gamma$ above is the one ArduPilot uses for a Quad-X frame, scaled by the Iris arm lengths.
\end{real}

% ---------------------------------------------------------------------
\subsection{All of Model 1 together: the closed loop}
\label{sec:m1loop}
% ---------------------------------------------------------------------

The last section of \code{M1\_3} connects everything as in \cref{fig:loop}: the true drone (block 1) $\rightarrow$
noisy gyro and camera $\rightarrow$ EKF (block 4) $\rightarrow$ controller on the \emph{estimate} (blocks 5--7)
$\rightarrow$ mixer (block 8) $\rightarrow$ the true drone, in gusty wind, following a climbing helix for 25\,s.
After the start-up transient the drone follows the helix with an RMS error of \textbf{4.7\,cm}, the estimator is
accurate to \textbf{2.2\,cm}, and the motors stay between 649 and 664\,rad/s.

\begin{figure}[H]
  \centering
  \includegraphics[width=0.6\linewidth]{M1_3_SO3_Control_Closed_Loop_2.png}
  \caption{The complete Model 1 loop in wind: desired helix (dotted), true flight (blue), EKF estimate (red, dashed).}
\end{figure}

\begin{toycode}{M1\_3\_SO3\_Control\_Closed\_Loop}{The closed loop}
ref = DroneLib.helix(t(k));                                   % mission
[ph, vh, Rh, wh] = DroneLib.unpack(xh);                       % the ESTIMATE
u = DroneLib.so3_control(ph, vh, Rh, wh, ref, P);             % blocks 5-7
[Om(:,k), ~, u_app] = DroneLib.mixer(u, P);                   % block 8
x  = DroneLib.rk4(x, u_app, dt, P, gust);                     % block 1, the real drone
[xh, Pk] = DroneLib.ekf_predict(xh, Pk, u_app, dt, P, Q);     % block 4
\end{toycode}

\subsection*{Recap of Model 1}
\begin{center}\small
\begin{tabularx}{\linewidth}{@{}c >{\raggedright\arraybackslash}p{2.8cm} >{\raggedright\arraybackslash}p{2.9cm} >{\raggedright\arraybackslash}p{2.4cm} X@{}}
\toprule
\textbf{\#} & \textbf{Block} & \textbf{In} & \textbf{Out} & \textbf{The idea in one sentence}\\
\midrule
1 & Dynamics & thrust, torque & motion & Push along the body's up-axis, gravity pulls down, torque spins the body.\\
2 & Linearization & operating point & $A$, $B$ & Zoomed in, the drone behaves like a linear system.\\
3 & Error dynamics & $A$, $B$, $\Delta t$ & $F$ & Small errors grow linearly, and fast ($t^3$), without feedback.\\
4 & EKF (VIO) & commands, sensors & $\hx_k$, $P_k$ & Predict with physics, correct with camera and IMU, weighted by trust.\\
5 & Position control & $\hp,\hv$, target & $\va_d$ & Pull towards the target like a spring with a damper.\\
6 & Force \& attitude & $\va_d$, yaw & $\vF_d$, $R_d$ & To accelerate sideways, lean; thrust direction = force direction.\\
7 & SO(3) attitude & $\hR,\hw,R_d,\vF_d$ & $\vtau$, $T$ & Twist downhill on the attitude error; the energy always decreases.\\
8 & Allocation & $T$, $\vtau$ & $\Omega_1..\Omega_4$ & Share thrust and torque among four rotors; speed $\propto\sqrt{\text{thrust}}$.\\
\bottomrule
\end{tabularx}
\end{center}
